\textbf{Mise en situation :} Après avoir survécu à votre session au Cégep, vous partez en vacances au village vacance lunaire. Alors que vous détendez vos muscles en apesanteur apparente dans le gymnase du complexe, un collègue vous défie de mesurer la masse d'un mystérieux haltère extraterrestre extrêmement dense. Vous n'avez évidemment pas pensé à apporter un pèse-personne dans vos bagages à main. Heureusement, votre trousse d'urgence contient un ressort parfait d'exercice physique !

\textbf{Problème :} Vous accrochez l'haltère au ressort dont la constante de rappel est $k = \v{Raideur}\text{ N/m}$. Vous étirez l'ensemble de $\v{Etirement}\text{ cm}$ depuis sa position d'équilibre, puis vous le relâchez sans vitesse initiale à $t = 0\text{ s}$. Vous observez que l'haltère repasse pour la toute première fois par sa position d'équilibre à l'instant $t = \v{TempsPassage}\text{ s}$.

\begin{enumerate}[label=\alph*)]
    \item Déterminez la période d'oscillation du système masse-ressort.
    \item Déterminez la masse de l'haltère mystérieux.
    \item Quelle est l'énergie mécanique totale emmagasinée dans le système ?
    \item Déterminez le module de la vitesse de l'haltère lorsqu'il atteint la position $x = \v{PosEval}\text{ m}$.
\end{enumerate}